\begin{thebibliography}{99}
\bibitem{r1n-ref1} Brown, T. B. et al. Language models are few-shot learners. In Advances in Neural Information Processing Systems 33 (eds Larochelle, H. et al.) (ACM, 2020). 题名中译: 语言模型是少样本学习器.
\bibitem{r1n-ref2} OpenAI et al. GPT4 technical report. Preprint at \url{https://doi.org/10.48550/arXiv.2303.08774} (2024). 题名中译: GPT4 技术报告.
\bibitem{r1n-ref3} Wei, J. et al. Chain-of-thought prompting elicits reasoning in large language models. In Advances in Neural Information Processing Systems 35 (eds Koyejo, S. et al.) 24824--24837 (ACM, 2022). 题名中译: 思维链提示激发大语言模型中的推理.
\bibitem{r1n-ref4} Wei, J. et al. Emergent abilities of large language models. In Transactions on Machine Learning Research (eds Kamath, G. et al.) (2022). 题名中译: 大语言模型的涌现能力.
\bibitem{r1n-ref5} Kaplan, J. et al. Scaling laws for neural language models. Preprint at \url{https://doi.org/10.48550/arXiv.2001.08361} (2020). 题名中译: 神经语言模型的缩放定律.
\bibitem{r1n-ref6} Kojima, T., Gu, S. S., Reid, M., Matsuo, Y. \& Iwasawa, Y. Large language models are zero-shot reasoners. In Advances in Neural Information Processing Systems 35 (eds Oh, A. H. et al.) 22199--22213 (ACM, 2022). 题名中译: 大语言模型是零样本推理者.
\bibitem{r1n-ref7} Chung, H. W. et al. Scaling instruction-finetuned language models. J. Mach. Learn. Res. 25, 1--53 (2024). 题名中译: 扩展经过指令微调的语言模型.
\bibitem{r1n-ref8} DeepSeek-AI et al. DeepSeek-V3 technical report. Preprint at \url{https://doi.org/10.48550/arXiv.2412.19437} (2025). 题名中译: DeepSeek-V3 技术报告.
\bibitem{r1n-ref9} Shao, Z. et al. DeepSeekMath: pushing the limits of mathematical reasoning in open language models. Preprint at \url{https://doi.org/10.48550/arXiv.2402.03300} (2024). 题名中译: DeepSeekMath: 推进开放语言模型数学推理的极限.
\bibitem{r1n-ref10} Wang, X. et al. Self-consistency improves chain of thought reasoning in language models. In 11th International Conference on Learning Representations (ICLR, 2023). 题名中译: 自一致性改进语言模型中的思维链推理.
\bibitem{r1n-ref11} Hendrycks, D. et al. Measuring massive multitask language understanding. In 9th International Conference on Learning Representations (ICLR, 2021). 题名中译: 衡量大规模多任务语言理解.
\bibitem{r1n-ref12} Gema, A. P. et al. Are we done with MMLU? In Proc. 2025 Conference of the Nations of the Americas Chapter of the Association for Computational Linguistics: Human Language Technologies (eds Chiruzzo, L. et al.) Vol. 1 (Long Papers), 5069--5096 (ACL, 2025). 题名中译: 我们已经完成 MMLU 的工作了吗?.
\bibitem{r1n-ref13} Wang, Y. et al. MMLU-Pro: a more robust and challenging multi-task language understanding benchmark. In Advances in Neural Information Processing Systems 37 (eds Globersons, A. et al.) 95266--95290 (ACM, 2024). 题名中译: MMLU-Pro: 更稳健、更具挑战性的多任务语言理解基准.
\bibitem{r1n-ref14} Dua, D. et al. DROP: a reading comprehension benchmark requiring discrete reasoning over paragraphs. In Proc. 2019 Conference of the North American Chapter of the Association for Computational Linguistics: Human Language Technologies Vol. 1 (Long and Short Papers) (eds Burstein, J. et al.) 2368--2378 (ACL, 2019). 题名中译: DROP: 需要对段落进行离散推理的阅读理解基准.
\bibitem{r1n-ref15} Huang, Y. et al. C-EVAL: a multi-level multi-discipline Chinese evaluation suite for foundation models. In Advances in Neural Information Processing Systems 36 (eds Oh, A. et al.) 62991--63010 (ACM, 2023). 题名中译: C-EVAL: 面向基础模型的多层次、多学科中文评估套件.
\bibitem{r1n-ref16} Zhou, J. et al. Instruction-following evaluation for large language models. Preprint at \url{https://doi.org/10.48550/arXiv.2311.07911} (2023). 题名中译: 大语言模型的指令遵循评估.
\bibitem{r1n-ref17} Krishna, S. et al. Fact, fetch, and reason: a unified evaluation of retrieval-augmented generation. In Proc. 2025 Conference of the Nations of the Americas Chapter of the Association for Computational Linguistics: Human Language Technologies Vol. 1 (Long Papers) 4745--4759 (ACL, 2025). 题名中译: 事实、获取与推理: 检索增强生成的统一评估.
\bibitem{r1n-ref18} Rein, D. et al. GPQA: a graduate-level Google-proof Q\&A benchmark. Preprint at \url{https://doi.org/10.48550/arXiv.2311.12022} (2023). 题名中译: GPQA: 研究生水平、无法轻易用 Google 检索作答的问答基准.
\bibitem{r1n-ref19} OpenAI. Introducing SimpleQA; \url{https://openai.com/index/introducing-simpleqa/} (2024). 题名中译: 介绍 SimpleQA.
\bibitem{r1n-ref20} He, Y. et al. Chinese SimpleQA: a Chinese factuality evaluation for large language models. In Proc. 63rd Annual Meeting of the Association for Computational Linguistics Vol. 1 (Long Papers), 19182--19208 (ACL, 2025). 题名中译: Chinese SimpleQA: 大语言模型的中文事实性评估.
\bibitem{r1n-ref21} Xu, L. et al. CLUE: a Chinese Language Understanding Evaluation benchmark. In Proc. 28th International Conference on Computational Linguistics (eds Scott, D. et al.) 4762--4772 (International Committee on Computational Linguistics, 2020). 题名中译: CLUE: 中文语言理解评估基准.
\bibitem{r1n-ref22} Dubois, Y., Galambosi, B., Liang, P. \& Hashimoto, T. B. Length-controlled AlpacaEval: a simple way to debias automatic evaluators. Preprint at \url{https://doi.org/10.48550/arXiv.2404.04475} (2025). 题名中译: 长度控制的 AlpacaEval: 消除自动评估器偏差的简单方法.
\bibitem{r1n-ref23} Li, T. et al. From crowdsourced data to high-quality benchmarks: Arena-Hard and BenchBuilder pipeline. Preprint at \url{https://doi.org/10.48550/arXiv.2406.11939} (2024). 题名中译: 从众包数据到高质量基准: Arena-Hard 与 BenchBuilder 流程.
\bibitem{r1n-ref24} OpenAI. Introducing SWE-bench verified; \url{https://openai.com/index/introducing-swe-bench-verified/} (2024). 题名中译: 介绍 SWE-bench Verified.
\bibitem{r1n-ref25} Aider. Aider LLM leaderboards; \url{https://aider.chat/docs/leaderboards/} (2024). 题名中译: Aider 大语言模型排行榜.
\bibitem{r1n-ref26} Jain, N. et al. LiveCodeBench: holistic and contamination free evaluation of large language models for code. In 13th International Conference on Learning Representations (ICLR, 2024). 题名中译: LiveCodeBench: 大语言模型代码能力的全面、无污染评估.
\bibitem{r1n-ref27} Mirzayanov, M. Codeforces; \url{https://codeforces.com/} (2025). 题名中译: Codeforces.
\bibitem{r1n-ref28} Chinese Mathematical Society (CMS). Chinese National High School Mathematics Olympiad; \url{https://www.cms.org.cn/Home/comp/comp/cid/12.html} (2024). 题名中译: 中国全国高中数学奥林匹克.
\bibitem{r1n-ref29} Mathematical Association of America. American Invitational Mathematics Examination; \url{https://maa.org/maa-invitational-competitions} (2024). 题名中译: 美国数学邀请赛.
\bibitem{r1n-ref30} OpenAI. Hello GPT-4o; \url{https://openai.com/index/hello-gpt-4o/} (2024). 题名中译: 你好, GPT-4o.
\bibitem{r1n-ref31} Schulman, J., Wolski, F., Dhariwal, P., Radford, A. \& Klimov, O. Proximal policy optimization algorithms. Preprint at \url{https://doi.org/10.48550/arXiv.1707.06347} (2017). 题名中译: 近端策略优化算法.
\bibitem{r1n-ref32} Ouyang, L. et al. Training language models to follow instructions with human feedback. In Advances in Neural Information Processing Systems 35 (eds Koyejo, S. et al.) 27730--27744 (ACM, 2022). 题名中译: 利用人类反馈训练遵循指令的语言模型.
\bibitem{r1n-ref33} Nano et al. deepseek-ai/DeepSeek-R1: v1.0.0. Zenodo \url{https://doi.org/10.5281/zenodo.15753192} (2025). 题名中译: deepseek-ai/DeepSeek-R1: v1.0.0.
\bibitem{r1n-ref34} Yu, X. et al. deepseek-ai/DeepSeek-V3: v1.0.0. Zenodo \url{https://doi.org/10.5281/zenodo.15753346} (2025). 题名中译: deepseek-ai/DeepSeek-V3: v1.0.0.
\bibitem{r1n-ref35} Paszke, A. et al. PyTorch: an imperative style, high-performance deep learning library. In Advances in Neural Information Processing Systems 32 (eds Wallach, H. M. et al.) 8026--8037 (ACM, 2019). 题名中译: PyTorch: 命令式、高性能的深度学习库.
\bibitem{r1n-ref36} Kwon, W. et al. Efficient memory management for large language model serving with PagedAttention. In Proc. ACM SIGOPS 29th Symposium on Operating Systems Principles 611--626 (ACM, 2023). 题名中译: 用 PagedAttention 高效管理大语言模型服务的内存.
\end{thebibliography}
